\section{Polymorphism}
\label{sec:polymorphism}

This section goes beyond the slides. In STLC and in typed LETREC, the identity
function has to be written once per type: \texttt{proc (x : int) x} and
\texttt{proc (x : bool) x} are different programs. Section~\ref{sec:inference}
showed that inference gives \texttt{proc (x) x} the type $t_0 \to t_0$, with a
variable left over. \textbf{Polymorphism} is the name for the features that let
one piece of code work at many types. There are three classical kinds:
\begin{itemize}
  \item parametric: one implementation that behaves the same at every type;
  \item ad hoc: one name with a different implementation per type
    (overloading, type classes);
  \item subtype: one implementation usable at every subtype of the type it
    expects (Section~\ref{sec:subtyping}).
\end{itemize}

\subsection{Parametric polymorphism}

System F abstracts over types the same way $\lambda$ abstracts over values. It
adds the type $\forall \alpha.\, \tau$, type abstraction $\Lambda \alpha.\, e$ and type
application $e\,[\tau]$:
\[
\mathit{id} = \Lambda \alpha.\, \lambda x{:}\alpha.\, x \;:\; \forall \alpha.\, \alpha \to \alpha.
\]
$\mathit{id}\,[\tyint]\ 3$ reduces to $3$ and $\mathit{id}\,[\tybool]\ \mathit{true}$ to
$\mathit{true}$: one term, instantiated twice (Figure~\ref{fig:instances}).
Hindley--Milner, the system behind ML and Haskell, is the fragment of System F
where $\forall$ only appears at the outside of a type and instantiation is
implicit. That restriction is what makes the inference of
Section~\ref{sec:inference} possible.

\begin{figure}[h]
\centering
\begin{forest}
  for tree={font=\small, s sep=10mm, l sep=8mm, edge={->}},
  [$\forall \alpha.\, \alpha \to \alpha$
    [$\tyint \to \tyint$, edge label={node[midway, left, font=\footnotesize]{$\alpha := \tyint$}}]
    [$\tybool \to \tybool$, edge label={node[midway, right, font=\footnotesize]{$\alpha := \tybool$}}]
  ]
\end{forest}
\caption{A polymorphic type and two of its instances.}
\label{fig:instances}
\end{figure}

\paragraph{Generalization happens at \texttt{let}.} In Haskell
\codelines{Haskell}{haskell/Polymorphism.hs}{3}{5}
type checks: the \texttt{let}-bound \texttt{ident} gets the type
$\forall \alpha.\, \alpha \to \alpha$, and each use instantiates it. The same thing with a
parameter, \lstinline|(\i -> (i 3, i True)) (\x -> x)|, is rejected with
\texttt{No instance for 'Num Bool'}: a parameter has one type in its whole
body, so after \lstinline|i 3| fixes it to a function on numbers, \lstinline|i True|
fails. Algorithm W quantifies the free type variables of a \texttt{let}-bound
value (they cannot appear in the context), and never those of a parameter.
ML restricts generalization further to values (the value restriction): a
generalized mutable cell could store a value at one type and be read at
another.

\subsection{Ad hoc polymorphism}

Overloading binds one name to unrelated implementations, chosen at compile
time by the types of the arguments:
\codelines{Java}{java/Overloading.java}{2}{3}
Haskell's type classes give this a structure: a class declares operations
and each instance implements them for one type. A function constrained by a
class is parametric in the type and ad hoc in the operation it calls:
\codelines{Haskell}{haskell/Polymorphism.hs}{7}{19}
Java gets the same effect with an interface as a bound on a type parameter
(Rust uses traits for this):
\codelines{Java}{java/Overloading.java}{5}{12}

\subsection{Higher-rank polymorphism}

Hindley--Milner only allows $\forall$ at the outside of a type. A function that
receives a polymorphic function and uses it at two types needs a $\forall$
inside the argument, a rank-2 type:
\codelines{Haskell}{haskell/Polymorphism.hs}{21}{22}
Inference cannot find this type (the first use would fix the type of
\texttt{f}), but checking against the written signature works. Haskell allows
it with the \texttt{RankNTypes} extension.

\subsection{Exercises}

\begin{exercise}
Write the System F term and type of \textit{const}, with
$\mathit{const}\ x\ y = x$. What is the normal form of
$\mathit{const}\,[\tyint]\,[\tybool]\ 3\ \mathit{true}$?
\end{exercise}

\paragraph{Solution.} $\Lambda \alpha.\, \Lambda \beta.\, \lambda x{:}\alpha.\, \lambda y{:}\beta.\, x : \forall \alpha.\, \forall \beta.\, \alpha \to \beta \to \alpha$.
The application reduces to $3$.

\begin{exercise}
What does the inference of Section~\ref{sec:inference} say about
\texttt{let i = proc (x) x in if (i zero?(0)) then (i 1) else 2}? Would the
program type check with let-polymorphism?
\end{exercise}

\paragraph{Solution.} It fails. The inference of Section~\ref{sec:inference} does
not generalize, so \texttt{i} has one type $t \to t$ in the whole body.
\texttt{(i zero?(0))} gives $t = \mathit{Boolean}$ and \texttt{(i 1)} gives
$t = \mathit{Int}$, and unification fails. With let-polymorphism \texttt{i} would get
$\forall t.\, t \to t$, each use would have its own instance, and the program
would have type $\mathit{Int}$.
